% ANCHOR: after_label def:c11-kernels
\begin{remark}[Physics notation: the composite line along $D$ and $P_m$]
The composite kernel is the line of the block after Section 9.2: the gauge propagator $\langle A_mA_n\rangle$, the classical scale transformation $(2x^n\partial_n+2)A_m$ of the gauge field inserted on it, and the line $\langle A^*_m\,c\rangle$, which in the gauge-fixed language is the ghost propagator joined to the antighost derivative $\partial_m\bar c$ that replaces the BRST source $A^*_m$ in Feynman gauge. Along a translation $P_m=\partial_m$ the inserted transformation is $b^n\partial_nA_m$. In momentum space the two propagators of the same momentum give the factor $e^{-2\varepsilon p^2}/p^4$ of the square kernel, so the line runs from a ghost $c$ to a gauge field $A_m$ and behaves as $|p|^{-3}$ at small momentum. It exists because $(Q^{GF})^2\neq0$ in Feynman gauge (block after Section 9.2), along the dilatation and along the translations alike; along the R-symmetry rotations it vanishes, the ghost and the gauge field at its two ends being R-singlets. The sector parts of $S^{(1)}(\Delta)$ and $S^{(1)}(P_b)$ are the scale transformation and the translation of the three field types separately: the scalars $A_I$, the gauginos $\lambda_\alpha{}^i,\bar\lambda_{\dot\alpha i}$, and the gauge field with the ghost.
\end{remark}

% ANCHOR: after_label def:c11-sp
\begin{remark}[Physics notation: regulator insertions and the composite pair kernel]
The special edge is the difference between the regularised and the bare contraction of a field with its BRST source, $e^{-\varepsilon p^2}-1$ times the field--source delta function: the regulator insertion by which the regularised BRST variation fails to be a derivation (block after Theorem \ref{thm:c4-obstruction}). A special-edge composite cycle is a one-loop graph through two constituent fields of a composite operator with one dilatation or translation insertion on the ghost--gauge line and one regulator insertion. The regulator insertion lives at proper times $t\le\varepsilon$, so the graph is a pure short-distance contribution: powers $\varepsilon^{-m}$, a finite part and no $\log\varepsilon$. The finite part is the composite pair kernel, a finite local two-body contribution of the gauge-fixing sector to the one-loop action of the dilatation and of the translations on composite operators; with the proper-time cutoff $\varepsilon^{-1/2}$ playing the role of the momentum cutoff (block after Proposition \ref{prop:c10-bound}), its value depends on the shape of the cutoff and is a scheme-dependent number.
\end{remark}

% ANCHOR: after_label thm:c11-local
\begin{remark}[Physics notation: locality of the gauge-fixing contribution]
In a Feynman-gauge computation of the one-loop scale transformation of a gauge-invariant composite operator, the graphs with the dilatation insertion on an internal gauge line split into graphs through two constituent fields and self-energy-type graphs on one field. The latter contain the composite two-point function, nonlocal on $A_m$, $\lambda$, $\bar\lambda$ and $c$ through $\log p^2$ of the field momentum (block after Theorem \ref{thm:c9-quadratic}). For a BRST-invariant operator the Slavnov--Taylor identity cancels these nonlocal terms against the graphs through pairs, by the argument in the proof of the theorem, and the whole insertion becomes the sum over pairs of constituent fields of the finite parts of the graphs with one regulator insertion: a local two-body operator, of the same form as the one-loop dilatation operator acting on neighbouring letters [B03c, (2.12)]. The one-body part of the one-loop dilatation is then the field-strength renormalisation of the letters together with the finite counterterm of the dilatation, as in the block after Theorem \ref{thm:c10-D1}.
\end{remark}

% ANCHOR: after_label prop:c11-tc0
\begin{remark}[Physics notation: the finite renormalisations of the scale transformation]
$\hbar F_1^{(1)}(\Delta)$ is the finite local counterterm of the dilatation. Up to BRST-exact field redefinitions its free quadratic part is a combination of the classical scale transformations $\varphi\mapsto(2x^m\partial_m+2w_\varphi)\varphi$ of the three field types, with the coefficients $\theta$ for the scalars $A_I$, $\theta_\lambda$ for the gauginos and $\theta_V$ for the gauge field with the ghost. Such a term changes the one-loop dilatation of a letter by a multiple of its classical scale transformation, a multiple of $x^m\partial_m+w$, and leaves the anomalous dimensions unchanged. Without reflection invariance the list has one further entry, the rotation $\lambda_\alpha{}^i\mapsto e^{i\varphi}\lambda_\alpha{}^i$, $\bar\lambda_{\dot\alpha i}\mapsto e^{-i\varphi}\bar\lambda_{\dot\alpha i}$ of the free gaugino, which is odd under the reflection of one direction of $V$ together with one direction of $W_{\mathbb R}^\perp$, the parity of Theorem \ref{thm:c9-quadratic} (iii).
\end{remark}

% ANCHOR: after_label thm:c11-arity2
\begin{remark}[Physics notation: the commutator of $D$ and $P_m$ at one loop]
With the physics dilatation $D=-E=-\tfrac12\Delta$ (key), the arity-two equation along $(\Delta,P_b)$ is the one-loop form of $[D,P_m]=P_m$ on composite operators, up to BRST-exact terms. In a scheme in which the translations act classically, it states that the part of the one-loop dilatation that is proportional to the position, $x^m$ times a one-loop translation, is BRST exact: the one-loop dilatation of an operator at the point $x$ is its anomalous-dimension operator, with no further term in $x^m\partial_m$. This condition fixes the coefficients $\theta$, $\theta_\lambda$, $\theta_V$ of the scale transformations in the finite counterterm, one for each field type, by requiring that the multiple of $x^m\partial_m$ produced by the gauge-fixing graphs cancel; it reproduces the values $-\tfrac13$ that the one-loop commutators of the regularised symmetry generators of the action fix (block after Section 9.4).
\end{remark}

% ANCHOR: after_label thm:c11-families
\begin{remark}[Physics notation: the light-cone families and the universal density]
$n^m$ with $n^2=0$ is a light-cone direction, $D_n=n^mD_m$ is the light-cone covariant derivative $D_+$, $\psi=\lambda_1{}^i$ is the $+$ component of a Weyl gaugino, and $F_{11}$ is the self-dual field-strength component $F_{++}$; the letters are the covariant letters $D_+^kZ$, $D_+^k\psi$, $D_+^kF_{++}$ of the block after Definition \ref{def:c10-letters}. The families are the light-cone $\mathfrak{sl}(2)$ sectors of [B03c]: the twist-two scalar operators $\operatorname{tr}(D_+^kZ\,D_+^{S-k}Z)$, the gaugino operators $\operatorname{tr}(D_+^k\psi\,D_+^{S-k}\psi)$ and the gluon operators $\operatorname{tr}(D_+^kF_{++}\,D_+^{S-k}F_{++})$, whose one-loop anomalous dimensions are the harmonic sums $\gamma=\frac{g^2N}{2\pi^2}h(j)$ with $j=S$, $S+1$, $S+2$ [B03c, (2.31)], together with the $\mathfrak{su}(1|1)$ sector of $Z$ and one gaugino component and two mixed families. The theorem says that, on two neighbouring fields of a single trace, the gauge-fixing contribution to the dilatation acts at the point of the operator as $-\tfrac43$ units times the identity on the same ordering, with no exchange and no trace term, and the contribution to the translation $P_m$ as $-\tfrac16$ units times the derivative of the pair: summed over a single trace, $-\tfrac43$ units for each field and $-\tfrac23$ units for each derivative, which is $-\tfrac23$ units times the letter count.
\end{remark}

% ANCHOR: after_label cor:c11-D1fam
\begin{remark}[Physics notation: anomalous dimensions and the absent one-loop translation]
The one-loop dilatation is $\tfrac12D_1$ (block after Definition \ref{def:c8-components}), so on the families $\tfrac12D_1=\gamma$ with $\gamma=\frac{g^2N}{8\pi^2}E$ and $E$ an eigenvalue of Beisert's one-loop Hamiltonian $\sum_lH_{l,l+1}$, $H_{12}=2h(J_{12})$: $\gamma=\frac{g^2N}{2\pi^2}h(S)$ on twist two, $\gamma=\frac{3g^2N}{4\pi^2}$ on the Konishi multiplet, which contains the trace scalar bilinear, $\operatorname{tr}(F_{++}F_{++})$ and the $\mathfrak{su}(1|1)$ state $\operatorname{tr}(Z\psi\psi)$, and $\frac{g^2N}{8\pi^2}$ times the spectra of the $\mathfrak{su}(1|1)$ table. For a one-loop quantisation through arity two, in the translation-covariant scheme $P_1$ vanishes on cohomology: the translations act on composite operators without a one-loop correction, $[P_m,O(x)]=\partial_mO(x)$, with the finite operator renormalisation of the block after Section 9.5 set to zero.
\end{remark}

% ANCHOR: after_label prop:c11-S
\begin{remark}[Physics notation: the special supercharges from the supercharges]
Classically the special supercharges $S^\alpha{}_i$, $\bar S^{\dot\alpha i}$ act on an operator at the point $x$ as the supercharges with the position-dependent parameter $\rho(x)\eta$, whose two-component form is $-i\,x_m\sigma^m_{\alpha\dot\alpha}\bar\eta^{\dot\alpha i}$ and $i\,x_m\eta^\alpha{}_i\sigma^m_{\alpha\dot\alpha}$, plus a term at the point: this is the content of the commutators $[P_m,S]\propto\sigma_m\bar Q$ and $[P_m,\bar S]\propto\sigma_mQ$ (Proposition \ref{prop:c2-special}). The proposition is the one-loop form of these commutators: when the one-loop translation vanishes, the position-dependent part of the one-loop special supercharge is the one-loop supercharge with the parameter $\rho(b)\eta$, and the one-loop special supercharge of an operator at the origin is the new datum. For the scalar bilinears, for which no such term at the point exists (block after Corollary \ref{cor:c7-S}), the one-loop $S$ is fixed by the one-loop $Q$. The one-loop $K_m$ is computed from its own regulated graphs at a finite upper proper-time endpoint, and $K=P^\dagger$ (Corollary \ref{cor:c8-gram}) relates it to the one-loop $P_m$ and the one-loop normalisation of the two-point functions as a check.
\end{remark}

% ANCHOR: after_label prop:c11-plane
\begin{remark}[Physics notation: the half-BPS bilinears of the twist]
For a rank-one supercharge, with parameter $\xi_\alpha{}^i=u_\alpha z^i$, such as the supercharge $Q^4_-$ of [BK] and [CY13, (2.1)] (C.-M. Chang and X. Yin, 1/16 BPS states in $\mathcal N=4$ super-Yang--Mills theory, arXiv:1305.6314), the scalars with $\rho(w)\xi=0$ are the three complex scalars $A_{ij}z^j$, the letters $\gamma^I\sim\Phi^{4I}$ of [BK, (3.51)], annihilated by $Q$ classically. Their symmetric traceless bilinears $\operatorname{tr}(\gamma^{(I}\gamma^{J)})$ are half-BPS chiral primaries in the $\mathbf{20'}$ (block after Section 9.4). The proposition states that their one-loop $Q$-variation vanishes identically, at every order at which the selection rules of the block after Proposition \ref{prop:c7-select} hold, which is the untwisted form of $Q_1(\gamma^I\gamma^J)=0$ [BK, (3.77)]: the statement holds for every finite renormalisation of the operators and of the symmetry generators and for every normalisation of the map from the fields to the twisted letters.
\end{remark}

% ANCHOR: end_of_section 11.6
\begin{remark}[Physics notation: the twisted one-loop supercharge and the Konishi entry]
The twisted differential $Q^\xi_0+\hbar Q^\xi_1$ is the action of a single Poincar\'e supercharge $Q$ of rank one on composite operators, classically and at one loop, and its cohomology at $\hbar^0$ is the space of classically $\tfrac1{16}$-BPS operators [CY13]. The one-loop part acts on pairs of letters [BK, (3.7)] and raises the dimension by $\tfrac12$, a two-site operator of order $g^2N$ in the canonical normalisation with the leg count of the order-$g^2$ two-site terms of the $\mathfrak{su}(2|3)$ supercharge of [B03] (N. Beisert, The su(2|3) dynamic spin chain, hep-th/0310252), with the leg count of the block after Section 10.9. The individual pair entries of [BK, (3.64)--(3.69)] depend on the normalisation of the operators: the Konishi entry, the one-loop term that maps the pair of a gaugino and a scalar to $\bar\lambda_{\dot\alpha}\bar\lambda^{\dot\alpha}$, is the twisted form of the one-loop Konishi anomaly, and its coefficient moves when the cubic scalar operator in the classical $Q$-variation is mixed with the operator $\operatorname{tr}(\bar\lambda_{\dot\alpha}\bar\lambda^{\dot\alpha})$ of the same quantum numbers. The scheme-independent object is the induced map $[Q^\xi_1]$ on the classical $Q$-cohomology (Theorem \ref{thm:c8-twist}), which decides which classically BPS operators acquire an anomalous dimension at one loop (block after Section 8.6).
\end{remark}

% ANCHOR: after_label lem:c11-lines
\begin{remark}[Physics notation: supersymmetry and the Feynman-gauge propagators]
The supercharges $Q^i_\alpha$, $\bar Q_{\dot\alpha i}$ commute with the free kinetic operators and with the heat kernel, $[\xi.,Q]=[\xi.,H]=0$, but not with the Feynman gauge fixing: the supersymmetry variation of the gauge field, $\delta_\xi A_m=i\big(\bar\lambda_{\dot\alpha i}\bar\sigma_m^{\dot\alpha\alpha}\xi_\alpha{}^i+\bar\xi_{\dot\alpha i}\bar\sigma_m^{\dot\alpha\alpha}\lambda_\alpha{}^i\big)$, changes the gauge-fixing function $\partial^mA_m$ by $\partial^m$ of a gaugino bilinear in $\xi$, and $C_\xi=[\xi.,Q^{GF}]$ is this change. On a loop, the insertion of $\delta_\xi A_m$ between the gaugino propagator and the propagator $\langle A^*_m\,c\rangle$, which in the gauge-fixed language is the ghost propagator joined to $\partial_m\bar c$, is a line from a gaugino to a ghost, $\mp2\,\xi\,e^{-2\varepsilon p^2}/p^2$ in momentum space: a scalar propagator carrying the supersymmetry parameter. The corresponding lines along the dilatation and the translations are annihilated by the BRST operator; this one is mapped to $(e^{-\varepsilon p^2}-e^{-Tp^2})$ times the variation of the gauge-fixing function, which is why the supersymmetry Ward identities of the regularised theory carry a ghost--gaugino term. The special supercharges $S^\alpha{}_i$, $\bar S^{\dot\alpha i}$ act as the supercharges with the position-dependent parameters $\bar\xi_{\dot\alpha i}=i\,x_m\eta^\alpha{}_i\sigma^m_{\alpha\dot\alpha}$, $\xi_\alpha{}^i=-i\,x_m\sigma^m_{\alpha\dot\alpha}\bar\eta^{\dot\alpha i}$, plus a term at the point that rotates a scalar into a gaugino, and the constant operator $Y_\eta$ that integration by parts produces; the special conformal transformations $K_m$ act on the Feynman-gauge propagator by the scale factor $4(b\cdot x)$ plus a constant term, the familiar failure of $\partial^mA_m$ to be conformally covariant.
\end{remark}

% ANCHOR: after_label thm:c11-classes
\begin{remark}[Physics notation: the one-loop superconformal action as a two-site operator]
For every local composite operator $O(x)$, BRST invariant or not, the one-loop correction to its transformation under $Q^i_\alpha$, $\bar Q_{\dot\alpha i}$, $S^\alpha{}_i$, $\bar S^{\dot\alpha i}$ and $K_m$ is a sum over pairs of constituent fields of a local two-site density, the finite part of the one-loop graphs through the two constituents with one insertion of the generator and one insertion of the regulator, plus a one-site term fixed by the finite counterterm of the symmetry, the analogue of a field renormalisation. In spin-chain language the one-loop supercharge acts on two neighbouring letters and raises the dimension by $\tfrac12$, like the two-site terms of order $g^2$ of the supercharge of the $\mathfrak{su}(2|3)$ dynamic chain [B03]; $S$ lowers it by $\tfrac12$ and $K_m$ by $1$. The regulator insertion lives at proper times below the cutoff, so the two-site densities are finite numbers without logarithms of the cutoff, and the cancellation of the poles and logarithms is the statement that the counterterms of the operators are compatible with the symmetry. A regulator insertion on a propagator that carries a divergent subgraph has a finite limit, which cancels against the same insertion on the counterterm of the subgraph; such finite counterterm insertions are therefore absent from the two-site densities. The special conformal transformations need the upper proper-time cutoff $T$, which plays the role of an infrared regulator of the special conformal Ward identity; the densities are independent of it.
\end{remark}

% ANCHOR: after_label thm:c11-bilinears
\begin{remark}[Physics notation: the one-loop supercharge of the $\mathbf{20'}$ and Konishi bilinears]
For $O=\operatorname{tr}(A_IA_J)w_1^Iw_2^J$, built from the six scalars $A_I$, the classical variation $Q^i_\alpha O$ replaces one scalar by a gaugino. On the twenty traceless bilinears, the half-BPS $\mathbf{20'}$, the one-loop variation is $c_Q$ times the classical one: a finite rescaling of the supercharge on these operators, which depends on the scheme through the finite counterterms and the relative normalisation of the gaugino and the scalar, and leaves the representation content unchanged. Its two-site part is $-2\log\tfrac32$ units, that is $\log\tfrac32\cdot\frac{g^2N}{16\pi^2}$ times the classical variation, the finite part of the graphs with the proper-time cutoff on every propagator, a scheme-dependent number; with the finite counterterm of the supercharge on the scalar leg the full coefficient is $\tfrac{79}{12}-3\log\tfrac32-t_Q$ units, and on the $\mathbf{20'}$ the combinations $c_Q+c_{\bar Q}-p_1$ and $c_Q+c_S-\tfrac14D_1$ equal $\tfrac{27}4-8\log\tfrac32$ units for every choice of the finite counterterms of the symmetry generators (Corollary \ref{cor:c11-F0}). The trace bilinear, the Konishi operator, receives in addition $4\log\tfrac32$ units times $(w_1,w_2)\operatorname{tr}(A_M\Gamma^M\xi\lambda)$ with the ten-dimensional connection $A_M=(A_m,A_I)$, including the gauge-variant term $\operatorname{tr}(A_m\bar\lambda\bar\sigma^m\xi)$. This term goes together with the one-loop BRST variation of the Konishi operator, proportional to $\operatorname{tr}(\partial_mc\,A_m)$: the finite counterterm $\operatorname{tr}(A_mA_m)$ added to the Konishi operator, which makes it a BRST cocycle in this gauge-fixed cutoff scheme, removes both, and the one-loop supersymmetry variation of the Konishi operator is then gauge invariant, as the closure of the BRST and supersymmetry transformations requires.
\end{remark}

% ANCHOR: after_label prop:c11-SKbil
\begin{remark}[Physics notation: $S$ and $K_m$ on the scalar bilinears]
On a scalar operator at the point $x$ the special supercharges act as the supercharges with the parameter $x_m\sigma^m\bar\eta$ or $x_m\eta\sigma^m$, so the two-site part of the one-loop $S$ on the bilinears is the one-loop $\bar Q$ or $Q$ evaluated at that parameter, with no term at the point. For a scalar primary of dimension $\Delta=2+\gamma$, $[K_b,O(x)]=\big(x^2\,b\cdot\partial-2(b\cdot x)\,x\cdot\partial-2\Delta\,(b\cdot x)\big)O(x)$, whose one-loop part is $-2\gamma(b\cdot x)O(x)$. The proposition reproduces this from the defect graphs: on the $\mathbf{20'}$, $\gamma=0$ and the one-loop $K$ vanishes, and on the Konishi operator, $\gamma=\frac{3g^2N}{4\pi^2}$, it is $-\frac{3g^2N}{2\pi^2}(b\cdot x)O$, which is $48$ units times $(b\cdot x)O$; at the origin the one-loop $K$ vanishes on both, as on every primary.
\end{remark}

% ANCHOR: after_label prop:c11-relations
\begin{remark}[Physics notation: the superconformal algebra at one loop]
The relations are the one-loop parts of the commutators of $\mathfrak{psl}(4|4)$ acting on gauge-invariant composite operators, valid up to BRST-exact terms: (R1) the one-loop supersymmetry variation of a BRST-invariant operator is BRST invariant; (R2) is $\{Q^i_\alpha,\bar Q_{\dot\beta j}\}\propto\delta^i_j\sigma^m_{\alpha\dot\beta}P_m$ and $\{Q,Q\}=0$; (R3) is $\{Q,S\}=$ dilatation plus Lorentz plus R-symmetry, which on the Konishi operator ties the one-loop $Q$ and $S$ to its anomalous dimension $\frac{3g^2N}{4\pi^2}$; (R4) is $\mathcal Q^2=0$ for the twisting supercharge at one loop; (R5) is $[K_m,P_n]$ and (R6) $\{S,\bar S\}\propto K$. On operators built from scalars alone, the extra terms vanish on the free equations of motion.
\end{remark}

% ANCHOR: after_label thm:c11-twistentries
\begin{remark}[Physics notation: the one-loop twisted supercharge on two letters]
For the supercharge $\mathcal Q=Q^4_-$ of [CY13] and [BK], the letters of the twisted theory are, by [BK, (3.51)], the three complex scalars $\gamma^I\sim\Phi^{4I}$, the gauginos $\beta_I\sim\Psi_{I+}$ and $\partial_{\dot\alpha}c\sim\bar\Psi^4_{\dot\alpha}$, the field strength $b\sim F_{++}$, and the holomorphic derivatives $\partial_{\dot\alpha}\sim D_{+\dot\alpha}$. [BK, (3.64)--(3.69)] give the one-loop $\mathcal Q_1$ on pairs of letters, each entry $\tfrac12f_{ACD}f_{BCE}$ times $\partial_{\dot\alpha}X^D\partial^{\dot\alpha}Y^E$ with equal colour words; the entry $\beta\gamma\to\partial c\,\partial c$ is the twisted form of the Konishi anomaly, the one-loop term proportional to $g^2N\operatorname{tr}\bar\lambda_{\dot\alpha}\bar\lambda^{\dot\alpha}$ in the supersymmetry variation of a Konishi descendant. In the untwisted regularised computation of this chapter, at zero finite renormalisation of the operators, the Konishi entry vanishes, the entries $\beta\beta$ and $\beta b\to\gamma\gamma$ have the structure of [BK], and $b\gamma$ and $\beta b\to\beta\,\partial c$ have unequal colour words and a $\log\tfrac32$. These numbers are scheme dependent: a finite mixing of the weight-$3$ operator $\operatorname{tr}(\Phi[\Phi,\Phi])$ with $\operatorname{tr}\bar\lambda\bar\lambda$ moves the Konishi entry, and the Konishi anomaly reappears in the class of $\mathcal Q_1$ on the $\mathcal Q$-cohomology. The quantity to compare with [BK] and with the lifting of $\tfrac1{16}$-BPS states is that class, which decides which classically BPS operators acquire an anomalous dimension at one loop.
\end{remark}
